\section{Related Work}
\label{sec:related}

\paragraph{Linguistic puzzles.} Olympiad puzzles probe rule induction from
small evidence sets \citep{lingoly,modeling,unveiling,majmudar2025};
obfuscation studies show part of measured performance is memorization
\citep{lingolytoo}. These works characterize behavior; we test whether the
inferred rules correspond to internal computation.

\paragraph{Program induction and verified reasoning.} Solvers have been
used to synthesize linguistic theories \citep{ellis2022} and to verify
LLM-proposed hypotheses or reasoning
\citep{wang2023hypothesis,parab2025,lyu2023,pan2023,olausson2023,ye2023}.
We repurpose verification: the certified program is not a way to produce an
answer but the specification of interpretability hypotheses.

\paragraph{Causal interpretability.} Probing establishes decodability but
not causal use \citep{conneau2018,hewitt2019,elazar2021,belinkov2022};
activation patching \citep{vig2020,meng2022,zhang2024} and DAS
\citep{geiger2021,geiger2024das,wu2023boundless,geiger2025} localize causal
variables, with path patching extending to circuits
\citep{wang2022ioi,goldowskydill2023,conmy2023}. The low-rank orthonormal
subspace parameterization we use is standard practice in this line, as
implemented in pyvene \citep{wu2024pyvene} and used by ReFT
\citep{wu2024reft} and AxBench \citep{wu2025axbench}. All of these methods
presuppose a known target variable; our targets come from task
certificates. Composition of features and interventions has been studied
via steering and feature geometry
\citep{turner2023,rimsky2024,merullo2024,bricken2023,huben2024,park2024};
certified rules let us test composition against operations the task
provably requires.
